\documentclass[a4paper]{article}
\usepackage{amsfonts}
\usepackage{amsmath}
\usepackage{amssymb}
\usepackage{graphicx}

\usepackage{cancel}

\begin{document}
  
\noindent \large Auteur: Siebe Mees \\
\noindent \large Studierichting: Informatica\\
\noindent \large Oefeningenreeks 6D nr. 31.21\\

\medskip

\normalsize

$ \displaystyle\sum_{n=0}^{\infty} \dfrac{2n+2}{\sqrt{n^3+5}} $\\

6+ Verhoudigncriterium \\

$ \lim\limits_{n \rightarrow \infty} \dfrac{\dfrac{2n+2}{\sqrt{n^3+5}}}{\dfrac{1}{\sqrt{n}}} = \lim\limits_{n \rightarrow \infty} \dfrac{2n\sqrt{n}+2\sqrt{n}}{\sqrt{n^3+5}} = \lim\limits_{n \rightarrow \infty} \dfrac{(2n+2)\sqrt{n}}{\sqrt{n^3+5}} \left(\sqrt{n^3+5} \sim n\sqrt{n} \right)$\\

$ = \lim\limits_{n \rightarrow \infty} \dfrac{2n\sqrt{n}}{n\sqrt{n}} = 2 \Rightarrow 0 < 2 < \infty $\\

$\Rightarrow \displaystyle\sum_{n=0}^{\infty} \dfrac{2n+2}{\sqrt{n^3+5}} \sim \displaystyle\sum_{n=1}^{\infty} \dfrac{1}{\sqrt{n}} $\\

$ \displaystyle\sum_{n=1}^{\infty} \dfrac{1}{\sqrt{n}} = \displaystyle\sum_{n=1}^{\infty} \dfrac{1}{n^{\tfrac{1}{2}}} $ is een $ p $-reeks met $ p=\dfrac{1}{2}<1 \Rightarrow $ divergent, dus $\displaystyle\sum_{n=0}^{\infty} \dfrac{2n+2}{\sqrt{n^3+5}} $ ook\\

\begin{thebibliography}{999}
%\bibitem{a} Referentie
\end{thebibliography}
\end{document} 
